\documentclass[12pt,a4paper]{article}

% ── Packages ──────────────────────────────────────────────────────────────────
\usepackage[T1]{fontenc}
\usepackage[utf8]{inputenc}
\usepackage[english]{babel}
\usepackage{lmodern}
\usepackage{geometry}
\usepackage{graphicx}
\usepackage{xcolor}
\usepackage{hyperref}
\usepackage{listings}
\usepackage{booktabs}
\usepackage{longtable}
\usepackage{array}
\usepackage{enumitem}
\usepackage{titlesec}
\usepackage{fancyhdr}
\usepackage{parskip}
\usepackage{microtype}
\usepackage{caption}
\usepackage{float}
\usepackage{amsmath}
\usepackage{mdframed}

% ── Page geometry ─────────────────────────────────────────────────────────────
\geometry{
  top=2.5cm, bottom=2.5cm,
  left=2.5cm, right=2.5cm,
  headheight=15pt
}

% ── Colours ───────────────────────────────────────────────────────────────────
\definecolor{f1red}{HTML}{E10600}
\definecolor{codebg}{HTML}{1F1F27}
\definecolor{codetext}{HTML}{E8E8F0}
\definecolor{codekw}{HTML}{FF7B72}
\definecolor{codestr}{HTML}{A5D6FF}
\definecolor{codecomment}{HTML}{8B949E}
\definecolor{codenumber}{HTML}{79C0FF}
\definecolor{lightgray}{HTML}{F5F5F5}
\definecolor{darkgray}{HTML}{555555}

% ── Hyperref setup ────────────────────────────────────────────────────────────
\hypersetup{
  colorlinks=true,
  linkcolor=darkgray,
  urlcolor=darkgray,
  citecolor=darkgray,
  pdftitle={F1 Knowledge Graph Explorer -- 1stWeb Semantics Report},
  pdfauthor={Rodrigo Abreu},
}

% ── Section styling ───────────────────────────────────────────────────────────
\titleformat{\section}
  {\large\bfseries}
  {\thesection.}{0.6em}{}
  [\vspace{-0.4em}\color{darkgray}\rule{\linewidth}{1.5pt}]

\titleformat{\subsection}
  {\normalsize\bfseries\color{darkgray}}
  {\thesubsection.}{0.5em}{}

\titleformat{\subsubsection}
  {\normalsize\itshape\color{darkgray}}
  {\thesubsubsection.}{0.5em}{}

% ── SPARQL listing style ──────────────────────────────────────────────────────
\lstdefinelanguage{SPARQL}{
  morekeywords={PREFIX,SELECT,WHERE,OPTIONAL,FILTER,GROUP,BY,ORDER,LIMIT,
                OFFSET,DISTINCT,COUNT,MIN,MAX,SUM,SAMPLE,INSERT,DELETE,DATA,
                ASC,DESC,UNION,BOUND,STR,CONTAINS,STRSTARTS,IF,IN,AS,FROM,
                GRAPH,CONSTRUCT,ASK,DESCRIBE,HAVING,VALUES,BIND,NOT,EXISTS},
  sensitive=true,
  morestring=[b]",
  morestring=[b]',
  morecomment=[l]{\#},
}

\lstdefinestyle{sparql}{
  language=SPARQL,
  backgroundcolor=\color{codebg},
  basicstyle=\ttfamily\footnotesize\color{codetext},
  keywordstyle=\color{codekw}\bfseries,
  stringstyle=\color{codestr},
  commentstyle=\color{codecomment}\itshape,
  numberstyle=\tiny\color{codenumber},
  breaklines=true,
  breakatwhitespace=true,
  tabsize=2,
  showstringspaces=false,
  frame=single,
  framesep=6pt,
  framerule=0pt,
  rulecolor=\color{codebg},
  xleftmargin=6pt,
  xrightmargin=6pt,
  captionpos=b,
}

\lstdefinestyle{shell}{
  backgroundcolor=\color{codebg},
  basicstyle=\ttfamily\footnotesize\color{codetext},
  breaklines=true,
  frame=single,
  framesep=6pt,
  framerule=0pt,
  rulecolor=\color{codebg},
  xleftmargin=6pt,
  xrightmargin=6pt,
}

\lstdefinestyle{python}{
  language=Python,
  backgroundcolor=\color{codebg},
  basicstyle=\ttfamily\footnotesize\color{codetext},
  keywordstyle=\color{codekw}\bfseries,
  stringstyle=\color{codestr},
  commentstyle=\color{codecomment}\itshape,
  breaklines=true,
  frame=single,
  framesep=6pt,
  framerule=0pt,
  rulecolor=\color{codebg},
  xleftmargin=6pt,
  xrightmargin=6pt,
  captionpos=b,
}

% ── Header / Footer ───────────────────────────────────────────────────────────
\pagestyle{fancy}
\fancyhf{}
\fancyhead[L]{\small\color{darkgray} F1 Knowledge Graph Explorer}
\fancyhead[R]{\small\color{darkgray} Web Semantics}
\fancyfoot[C]{\small\thepage}
\renewcommand{\headrulewidth}{0.4pt}
\renewcommand{\footrulewidth}{0pt}
\renewcommand{\headrule}{\color{darkgray}\rule{\linewidth}{0.4pt}}

% ── Info box ──────────────────────────────────────────────────────────────────
\newmdenv[
  backgroundcolor=lightgray,
  linecolor=f1red,
  linewidth=3pt,
  leftline=true, rightline=false, topline=false, bottomline=false,
  innerleftmargin=10pt,
  innertopmargin=8pt,
  innerbottommargin=8pt,
]{infobox}

% ─────────────────────────────────────────────────────────────────────────────
\begin{document}
% ─────────────────────────────────────────────────────────────────────────────

% ── Title page ────────────────────────────────────────────────────────────────
\begin{titlepage}
  \centering
  \vspace*{1.5cm}

  {\fontsize{48}{52}\selectfont\bfseries\color{f1red} F\textcolor{black}{1}}\\[0.4cm]
  {\Large\bfseries Knowledge Graph Explorer}\\[0.3cm]
  {\large\color{darkgray} Web Semantics --- Practical Assignment Report}

  \vspace{1.5cm}
  \color{darkgray}\rule{0.4\linewidth}{2pt}\\[0.1cm]
  \color{darkgray}\rule{\linewidth}{0.4pt}
  \vspace{1cm}

  \begin{tabular}{rl}
    \textbf{Author}      & Rodrigo Abreu (113626), Hugo Ribeiro (113402) \\[0.3em]
    \textbf{Course}      & Web Semantics \\[0.3em]
    \textbf{Year}        & 2025--2026 \\[0.3em]
    \textbf{Technology}  & Django 4 $\cdot$ GraphDB $\cdot$ SPARQL $\cdot$ RDF \\
  \end{tabular}

  \vspace{1.5cm}
  \color{darkgray}\rule{\linewidth}{0.4pt}\\[0.1cm]
  \color{darkgray}\rule{0.4\linewidth}{2pt}

  \vfill
  {\small\color{darkgray}\today}
\end{titlepage}

% ── Table of contents ─────────────────────────────────────────────────────────
\tableofcontents
\newpage

% =============================================================================
\section{Introduction}
% =============================================================================

Formula 1 is the premier class of single-seater auto racing sanctioned by the
Fédération Internationale de l'Automobile (FIA). Since its inaugural season in
1950, it has grown into one of the most data-rich sports in the world, with
precise records of every driver, constructor, race result, qualifying session,
lap time, and pit stop stored in structured datasets.

This project constructs a \textbf{semantic web application} around Formula 1
historical data covering the 1950--2024 seasons. The data is modelled as an RDF
knowledge graph stored in \textbf{GraphDB} (a native RDF triple store with
RDFS+ inference), exposed through a \textbf{SPARQL} query interface, and
presented to users via a \textbf{Django} web application styled with the visual
identity of the sport.

\subsection{Objectives}

\begin{itemize}[leftmargin=1.5em]
  \item Transform flat CSV data from the Formula 1 World Championship (1950 - 2024) dataset into
        a well-structured RDF knowledge graph using custom Python scripts.
  \item Model the domain with a consistent namespace and predicate vocabulary
        that enables efficient SPARQL queries over millions of triples.
  \item Build a full-featured web UI allowing public users to explore drivers,
        constructors, circuits, races, and seasons in detail.
  \item Provide an authenticated admin panel allowing staff to create, update,
        and delete core entities directly in the knowledge graph via SPARQL
        UPDATE, including relation-backed fields.
  \item Add higher-level knowledge-graph operations such as batch CSV import,
        dry-run preview, rollback, and natural-language querying over SPARQL.
\end{itemize}

\subsection{Technology Stack}

\begin{center}
\begin{tabular}{@{}lll@{}}
\toprule
\textbf{Layer}     & \textbf{Technology}      & \textbf{Role} \\
\midrule
Knowledge Store    & GraphDB 10 (Ontotext)    & RDF triple store, SPARQL endpoint \\
Query Language     & SPARQL 1.1               & SELECT, UPDATE (INSERT/DELETE) \\
Data Format        & N-Triples (.nt)          & RDF serialisation for bulk import \\
Web Framework      & Django 4.x (Python)      & Views, templates, URL routing \\
SPARQL Client      & SPARQLWrapper            & HTTP bridge to GraphDB endpoint \\
Serialisation      & RDFLib                   & CSV $\rightarrow$ RDF conversion \\
Auth               & Django contrib.auth      & Session-based login for admin \\
Frontend           & Vanilla HTML/CSS/JS      & Responsive dark-theme UI \\
LLM Integration    & Gemini API               & Natural-language to SPARQL assistant \\
\bottomrule
\end{tabular}
\end{center}

% =============================================================================
\section{Data, Sources, and Transformation Process}
% =============================================================================

\subsection{Data Source}

The raw data originates from the \textbf{Formula 1 World Championship (1950 - 2024)}
(\url{https://www.kaggle.com/datasets/rohanrao/formula-1-world-championship-1950-2020?select=circuits.csv}), which contains
structured Formula 1 statistics from 1950 to 2024, that can be downloaded as csv files. 

The dataset consists of \textbf{13 CSV files}, each representing a distinct
entity type or relationship:

\begin{center}
\begin{longtable}{@{}llp{6.5cm}@{}}
\toprule
\textbf{CSV File}              & \textbf{RDF Class}       & \textbf{Key Fields} \\
\midrule
\endhead
\midrule
\multicolumn{3}{r}{\small continued \ldots}\\
\endfoot
\bottomrule
\endlastfoot
\texttt{seasons.csv}           & Season                   & year, url \\
\texttt{circuits.csv}          & Circuit                  & circuitId, name, location, country, lat, lng, alt, url \\
\texttt{constructors.csv}      & Constructor              & constructorId, name, nationality, url \\
\texttt{drivers.csv}           & Driver                   & driverId, forename, surname, code, number, dob, nationality, url \\
\texttt{races.csv}             & Race                     & raceId, year, round, name, date, circuitId (FK), url \\
\texttt{results.csv}           & Result                   & resultId, raceId, driverId, constructorId, grid, positionOrder, points, laps, time, statusId, fastestLapTime, fastestLapSpeed \\
\texttt{qualifying.csv}        & QualifyingResult         & qualifyId, raceId, driverId, constructorId, position, q1, q2, q3 \\
\texttt{pit\_stops.csv}         & PitStop                  & raceId, driverId, stop, lap, milliseconds, duration \\
\texttt{lap\_times.csv}         & LapTime                  & raceId, driverId, lap, position, time, milliseconds \\
\texttt{driver\_standings.csv}  & DriverStanding           & driverStandingsId, raceId, driverId, points, position, wins \\
\texttt{constructor\_standings.csv} & ConstructorStanding & constructorStandingsId, raceId, constructorId, points, position, wins \\
\texttt{constructor\_results.csv}   & ConstructorResult   & constructorResultsId, raceId, constructorId, points, status \\
\texttt{status.csv}            & Status                   & statusId, status (text label) \\
\end{longtable}
\end{center}

\subsection{RDF Namespace and Vocabulary}

A custom namespace \texttt{http://example.org/f1/} is used for all domain
predicates, and \texttt{http://example.org/resource/} is used as the base for
all entity URIs. Standard W3C vocabularies are used where appropriate.

\begin{lstlisting}[style=sparql, caption={SPARQL namespace prefixes used throughout the application}]
PREFIX f1:   <http://example.org/f1/>
PREFIX res:  <http://example.org/resource/>
PREFIX rdf:  <http://www.w3.org/1999/02/22-rdf-syntax-ns#>
PREFIX rdfs: <http://www.w3.org/2000/01/rdf-schema#>
PREFIX xsd:  <http://www.w3.org/2001/XMLSchema#>
\end{lstlisting}

\subsubsection{URI Construction}

Entity URIs follow a consistent pattern:
\texttt{res:\{kind\}/\{id\}}, where \texttt{kind} is a lowercase slug derived
from the entity class and \texttt{id} is the numeric identifier from the source
CSV.

\begin{center}
\begin{tabular}{@{}ll@{}}
\toprule
\textbf{Entity}    & \textbf{Example URI} \\
\midrule
Driver             & \texttt{res:driver/1} \\
Constructor        & \texttt{res:constructor/6} \\
Circuit            & \texttt{res:circuit/6} \\
Race               & \texttt{res:race/841} \\
Result             & \texttt{res:result/1} \\
QualifyingResult   & \texttt{res:qualifying-result/1} \\
PitStop            & \texttt{res:pit-stop/\{raceId\}/\{driverId\}/\{stop\}} \\
LapTime            & \texttt{res:lap-time/\{raceId\}/\{driverId\}/\{lap\}} \\
Status             & \texttt{res:status/1} \\
\bottomrule
\end{tabular}
\end{center}

\subsection{Transformation Script}

The Python script \texttt{scripts/csv\_to\_rdf.py} performs the full
CSV-to-RDF conversion using RDFLib. The process for each CSV file is:

\begin{enumerate}[leftmargin=2em]
  \item Read each row with \texttt{csv.DictReader}.
  \item Construct the subject URI from the configured \texttt{id\_columns}.
  \item For each column, cast the value to the appropriate XSD datatype
        (\texttt{xsd:integer}, \texttt{xsd:decimal}, \texttt{xsd:date},
        \texttt{xsd:gYear}) or leave it as a plain literal.
  \item For foreign key columns (e.g.\ \texttt{raceId}, \texttt{driverId}),
        emit an object property triple linking to the target entity URI.
  \item Emit a \texttt{rdfs:label} triple with a human-readable label derived
        from available name fields.
  \item Serialise the complete graph as \textbf{N-Triples} (\texttt{.nt}) for
        efficient bulk import into GraphDB.
\end{enumerate}

\subsubsection{Data Type Mapping}

\begin{center}
\begin{tabular}{@{}lll@{}}
\toprule
\textbf{Column group}      & \textbf{XSD type}     & \textbf{Example columns} \\
\midrule
Numeric identifiers/counts & \texttt{xsd:integer}  & driverId, raceId, laps, position \\
Floating-point values      & \texttt{xsd:decimal}  & lat, lng, points, fastestLapSpeed \\
Calendar dates             & \texttt{xsd:date}     & dob, date, fp1\_date \\
Calendar years             & \texttt{xsd:gYear}    & year \\
Time strings               & plain literal         & time, fastestLapTime, q1, q2, q3 \\
External URLs              & \texttt{rdfs:seeAlso} & url \\
\bottomrule
\end{tabular}
\end{center}

\subsubsection{Derived Labels}

Because SPARQL queries use \texttt{rdfs:label} for display, the script
generates a label for every entity at conversion time. For drivers, the label
is the concatenation of \texttt{forename} and \texttt{surname}. For races it
uses the \texttt{name} column. For seasons the label is ``Season YYYY''.

\subsubsection{Sprint Results Exclusion}

The original dataset includes a \texttt{sprint\_results.csv} file recording
results from Saturday sprint races. Sprint results use the same
\texttt{resultId} sequence as regular race results and link to the same race
entities via \texttt{f1:race}. This caused \textbf{duplicate rows} when
querying race results: for any sprint weekend, both a \texttt{res:result/N}
and a \texttt{res:sprint-result/N} entity would match the \texttt{f1:resultId}
anchor, effectively doubling every driver's entry in the results table.

Sprint results were therefore excluded from the conversion script. A SPARQL
\texttt{FILTER} guard is also applied in the race detail query as a defensive
measure for any existing data in the store:

\begin{lstlisting}[style=sparql, caption={Filter excluding sprint-result URIs from race result queries}]
FILTER(!CONTAINS(STR(?res), "/sprint-result/"))
\end{lstlisting}

\subsection{Loading into GraphDB}

The generated \texttt{.nt} file is loaded into GraphDB using the \textbf{Import
$\rightarrow$ Upload RDF files} interface in the GraphDB Workbench, targeting
the named graph \texttt{http://example.org/graph/formula1}. The repository is
configured with \textbf{RDFS+ (Optimized)} inference enabled, which allows
entailments over \texttt{rdfs:subClassOf} and \texttt{rdfs:subPropertyOf}
chains.

% =============================================================================
\section{Operations on the Data (SPARQL Queries)}
% =============================================================================

All SPARQL queries are executed via the \texttt{GraphDBClient} class in
\texttt{championship/services/graphdb.py}, which prepends the five namespace
prefixes to every query. The client exposes two methods:

\begin{itemize}[leftmargin=1.5em]
  \item \texttt{query(sparql)} --- executes a SELECT and returns a list of
        dictionaries mapping variable names to string values.
  \item \texttt{run\_update(sparql)} --- executes a SPARQL UPDATE (INSERT or
        DELETE) against the \texttt{/statements} endpoint.
\end{itemize}

\subsection{Anchoring Strategy}

A central design decision is the use of \textbf{property anchors} rather than
\texttt{rdf:type} assertions for entity discovery. GraphDB's RDFS+ inference
engine processes \texttt{rdf:type} through the class hierarchy, which is
expensive on large graphs. Instead, every query starts with a property that is
unique to the target entity class:

\begin{center}
\begin{tabular}{@{}ll@{}}
\toprule
\textbf{Entity}    & \textbf{Anchor predicate} \\
\midrule
Driver             & \texttt{f1:driverId} \\
Constructor        & \texttt{f1:constructorId} \\
Circuit            & \texttt{f1:circuitRef} (not \texttt{f1:circuitId} --- races also carry this) \\
Race               & \texttt{f1:round} \\
Result             & \texttt{f1:resultId} \\
QualifyingResult   & \texttt{f1:qualifyId} \\
PitStop            & \texttt{f1:stop} \\
LapTime            & \texttt{f1:position} (lap\_times only; pit\_stops do not have this) \\
DriverStanding     & \texttt{f1:driverStandingsId} \\
ConstructorStanding & \texttt{f1:constructorStandingsId} \\
\bottomrule
\end{tabular}
\end{center}

\subsection{Year Filtering}

Race years are stored as \texttt{xsd:gYear}. Direct numeric comparison
(\texttt{FILTER(?year = 2024)}) fails because the literal type does not match
\texttt{xsd:integer}. All season-scoped queries therefore coerce to string:

\begin{lstlisting}[style=sparql]
FILTER(STR(?yr) = "2024")
\end{lstlisting}

\subsection{Read Queries (SELECT)}

\subsubsection{Homepage Statistics}

The home page displays live aggregate counts queried directly from the graph,
ensuring they always reflect the current state of the knowledge base.

\begin{lstlisting}[style=sparql, caption={Aggregate entity counts for the homepage}]
SELECT
  (COUNT(DISTINCT ?season)      AS ?seasons)
  (COUNT(DISTINCT ?circuit)     AS ?circuits)
  (COUNT(DISTINCT ?driver)      AS ?drivers)
  (COUNT(DISTINCT ?constructor) AS ?constructors)
  (COUNT(DISTINCT ?race)        AS ?races)
WHERE {
  ?race     f1:round       ?round ;
            f1:year        ?season ;
            f1:circuit     ?circuit .
  ?res      f1:resultId    ?anyId ;
            f1:race        ?race ;
            f1:driver      ?driver ;
            f1:constructor ?constructor .
}
\end{lstlisting}

\subsubsection{Driver List with Career Statistics}

Three queries are executed in sequence to build the drivers list page: one for
basic fields, one for career year ranges, and one for podium breakdowns.

\begin{lstlisting}[style=sparql, caption={Career year range per driver (MIN/MAX over race years)}]
SELECT ?driver
       (COUNT(*) AS ?races)
       (MIN(STR(?yr)) AS ?firstYear)
       (MAX(STR(?yr)) AS ?lastYear)
WHERE {
  ?r f1:resultId ?anyId ;
     f1:driver   ?driver ;
     f1:race     ?race .
  ?race f1:year ?yr .
}
GROUP BY ?driver
\end{lstlisting}

\begin{lstlisting}[style=sparql, caption={Podium breakdown (positions 1, 2, 3) per driver}]
SELECT ?driver ?pos (COUNT(*) AS ?cnt) WHERE {
  ?r f1:resultId    ?anyId ;
     f1:driver      ?driver ;
     f1:positionOrder ?pos .
  FILTER(?pos IN (1, 2, 3))
}
GROUP BY ?driver ?pos
\end{lstlisting}

\subsubsection{Constructor Detail --- Driver Roster}

This query uses conditional aggregation (\texttt{SUM(IF(\ldots))}) to compute
per-driver wins and podiums without leaving the SPARQL layer.

\begin{lstlisting}[style=sparql, caption={Driver roster with win and podium counts for a constructor}]
SELECT ?driver ?driverLabel
       (COUNT(*) AS ?races)
       (SUM(IF(?pos = 1, 1, 0))         AS ?wins)
       (SUM(IF(?pos IN (1,2,3), 1, 0))  AS ?podiums)
       (MIN(?year) AS ?firstYear)
       (MAX(?year) AS ?lastYear)
WHERE {
  ?r f1:resultId    ?anyId ;
     f1:constructor <{constructorUri}> ;
     f1:driver      ?driver ;
     f1:positionOrder ?pos ;
     f1:race        ?race .
  ?race   f1:year    ?year .
  ?driver rdfs:label ?driverLabel .
}
GROUP BY ?driver ?driverLabel
ORDER BY DESC(?wins) DESC(?races)
\end{lstlisting}

\subsubsection{Race Detail --- Results with Finish Status}

Results for a race are retrieved in a single query. The finish status (e.g.\
``Finished'', ``Engine'', ``Collision'') is joined via \texttt{f1:statusId}
linking results to status entities. Both joins are wrapped in \texttt{OPTIONAL}
so drivers without a recorded status still appear.

\begin{lstlisting}[style=sparql, caption={Race results with finish status join}]
SELECT ?positionOrder ?grid ?laps ?time ?points ?statusText
       ?driverLabel ?driverId ?constructorLabel WHERE {
  ?res f1:resultId     ?anyId ;
       f1:race         <{raceUri}> ;
       f1:positionOrder ?positionOrder ;
       f1:driver        ?driver ;
       f1:constructor   ?constructor .
  FILTER(!CONTAINS(STR(?res), "/sprint-result/"))
  OPTIONAL { ?res f1:grid   ?grid   }
  OPTIONAL { ?res f1:laps   ?laps   }
  OPTIONAL { ?res f1:time   ?time   }
  OPTIONAL { ?res f1:points ?points }
  OPTIONAL {
    ?res       f1:statusId ?statusId .
    ?statusEnt f1:statusId ?statusId ;
               f1:status   ?statusText .
    FILTER(!CONTAINS(STR(?statusEnt), "/result/"))
  }
  ?driver      rdfs:label ?driverLabel ;
               f1:driverId ?driverId .
  ?constructor rdfs:label ?constructorLabel .
}
ORDER BY ?positionOrder
\end{lstlisting}

\subsubsection{Pit Stop Aggregation}

Pit stop data is aggregated per driver per race. The \texttt{f1:stop} property
serves as the anchor (unique to pit stop entities). The query computes total pit
time, fastest single stop time, and the lap on which the fastest stop occurred.

\begin{lstlisting}[style=sparql, caption={Per-driver pit stop aggregation for a race}]
SELECT ?driverLabel ?driverId
       (COUNT(*)         AS ?stops)
       (SUM(?ms)         AS ?totalMs)
       (MIN(?ms)         AS ?fastestMs)
       (SAMPLE(?fastestLap) AS ?fastestLap)
WHERE {
  ?ps f1:stop         ?stop ;
      f1:race         <{raceUri}> ;
      f1:driver       ?driver ;
      f1:milliseconds ?ms ;
      f1:lap          ?fastestLap .
  ?driver rdfs:label ?driverLabel ;
          f1:driverId ?driverId .
}
GROUP BY ?driverLabel ?driverId
ORDER BY ?driverLabel
\end{lstlisting}

\subsubsection{Season Detail --- Championship Standings}

Standings are stored per-race (cumulative totals after each round). To obtain
\emph{final} standings for a season, the query orders by round descending and
the application keeps only the first occurrence per driver/constructor.

\begin{lstlisting}[style=sparql, caption={Driver championship standings for a season (final round)}]
SELECT ?driverLabel ?driverId
       ?points ?position ?wins ?round WHERE {
  ?ds f1:driverStandingsId ?anyId ;
      f1:race     ?race ;
      f1:driver   ?driver ;
      f1:points   ?points ;
      f1:position ?position ;
      f1:wins     ?wins .
  ?race f1:year  ?yr ;
        f1:round ?round .
  FILTER(STR(?yr) = "{year}")
  ?driver rdfs:label ?driverLabel ;
          f1:driverId ?driverId .
}
ORDER BY DESC(?round) ?position
\end{lstlisting}

\subsubsection{Circuit Fastest Lap}

Lap time entities carry a \texttt{f1:position} property (finishing position
within that lap). Position~1 identifies the fastest lap in each lap. However,
since the circuit detail page seeks the all-time fastest lap recorded at that
circuit across all races, the query retrieves the minimum milliseconds value.

\begin{lstlisting}[style=sparql, caption={Fastest laps ever recorded at a circuit}]
SELECT ?driverLabel ?driverId ?time ?ms ?year WHERE {
  ?race f1:circuit <{circuitUri}> ;
        f1:year     ?year .
  ?lt   f1:race         ?race ;
        f1:driver       ?driver ;
        f1:milliseconds ?ms ;
        f1:time         ?time ;
        f1:lap          ?lap ;
        f1:position     ?pos .
  ?driver rdfs:label ?driverLabel ;
          f1:driverId ?driverId .
  FILTER(?ms > 60000)
}
ORDER BY ASC(?ms)
LIMIT 5
\end{lstlisting}

\subsection{Write Operations (SPARQL UPDATE)}

The admin panel performs SPARQL UPDATE queries via the
\texttt{/statements} endpoint. All mutations are expressed using standard
SPARQL 1.1 UPDATE syntax.

\subsubsection{Entity Creation (INSERT DATA)}

New entities are created with a single \texttt{INSERT DATA} block. Before
inserting, a helper query computes the next available ID by finding the current
maximum for the target class:

\begin{lstlisting}[style=sparql, caption={Maximum ID query for sequential ID generation (shown for drivers; analogous for all entity types)}]
SELECT (MAX(?id) AS ?maxId) WHERE {
  ?e f1:driverId ?id .
}
\end{lstlisting}

The same pattern is applied for \texttt{f1:constructorId}, \texttt{f1:circuitId},
\texttt{f1:raceId}, etc.\ Each entity type has a well-defined set of required
and optional predicates:

\begin{lstlisting}[style=sparql, caption={INSERT DATA for a new driver with all fields}]
INSERT DATA {
  <http://example.org/resource/driver/862>
    rdfs:label     "George Russell" ;
    f1:driverId    862 ;
    f1:driverRef   "russell" ;
    f1:forename    "George" ;
    f1:surname     "Russell" ;
    f1:code        "RUS" ;
    f1:number      63 ;
    f1:dob         "1998-02-15"^^xsd:date ;
    f1:nationality "British" ;
    rdfs:seeAlso   <https://en.wikipedia.org/wiki/George_Russell> .
}
\end{lstlisting}

\begin{lstlisting}[style=sparql, caption={INSERT DATA for a new constructor}]
INSERT DATA {
  <http://example.org/resource/constructor/212>
    rdfs:label        "Aston Martin" ;
    f1:constructorId  212 ;
    f1:constructorRef "aston_martin" ;
    f1:name           "Aston Martin" ;
    f1:nationality    "British" ;
    rdfs:seeAlso      <https://en.wikipedia.org/wiki/Aston_Martin_in_Formula_One> .
}
\end{lstlisting}

\begin{lstlisting}[style=sparql, caption={INSERT DATA for a new circuit}]
INSERT DATA {
  <http://example.org/resource/circuit/78>
    rdfs:label    "Las Vegas Street Circuit" ;
    f1:circuitId  78 ;
    f1:circuitRef "las_vegas" ;
    f1:name       "Las Vegas Street Circuit" ;
    f1:location   "Las Vegas" ;
    f1:country    "USA" ;
    f1:lat        36.1699 ;
    f1:lng        -115.1398 ;
    rdfs:seeAlso  <https://en.wikipedia.org/wiki/Las_Vegas_Grand_Prix> .
}
\end{lstlisting}

\begin{lstlisting}[style=sparql, caption={INSERT DATA for a new race, linking to an existing circuit}]
INSERT DATA {
  <http://example.org/resource/race/1126>
    rdfs:label  "Las Vegas Grand Prix" ;
    f1:raceId   1126 ;
    f1:name     "Las Vegas Grand Prix" ;
    f1:year     "2025"^^xsd:gYear ;
    f1:round    22 ;
    f1:circuit  <http://example.org/resource/circuit/78> ;
    f1:date     "2025-11-22"^^xsd:date ;
    rdfs:seeAlso <https://en.wikipedia.org/wiki/2025_Las_Vegas_Grand_Prix> .
}
\end{lstlisting}

\begin{lstlisting}[style=sparql, caption={INSERT DATA for a new season}]
INSERT DATA {
  <http://example.org/resource/season/2025>
    f1:seasonYear "2025"^^xsd:gYear ;
    rdfs:label    "Season 2025" ;
    rdfs:seeAlso  <https://en.wikipedia.org/wiki/2025_Formula_One_World_Championship> .
}
\end{lstlisting}

\subsubsection{Entity Update (DELETE + INSERT)}

Updating an entity requires removing all its current triples first, then
reinserting the complete set of new values. Both operations are sent in a
single SPARQL UPDATE request, separated by a semicolon, which provides
pseudo-atomic semantics:

\begin{lstlisting}[style=sparql, caption={DELETE + INSERT pattern for entity updates (identical structure for all entity types)}]
DELETE WHERE {
  <http://example.org/resource/driver/862> ?p ?o
} ;
INSERT DATA {
  <http://example.org/resource/driver/862>
    rdfs:label     "George Russell" ;
    f1:driverId    862 ;
    f1:nationality "British" ;
    f1:number      63 ;
    ... .
}
\end{lstlisting}

This pattern is identical for all five entity types (drivers, constructors,
circuits, races, seasons). The only difference is the subject URI prefix and
the set of predicates emitted. Because the DELETE removes \emph{all} predicates
before the INSERT, optional fields that were previously set and are now absent
(e.g.\ a driver code that was cleared) are cleanly removed from the graph.

For relation-backed forms the same pattern is reused. For example, editing a
driver can now also replace or clear the \texttt{f1:constructor} object
property, allowing the admin panel to modify not only scalar literals but also
links between RDF entities.

\subsubsection{Entity Deletion (DELETE WHERE)}

Deleting an entity removes all triples where the entity appears as the subject:

\begin{lstlisting}[style=sparql, caption={DELETE WHERE for entity removal}]
DELETE WHERE {
  <http://example.org/resource/driver/862> ?p ?o .
}
\end{lstlisting}

The same one-liner applies to constructors, circuits, races, and seasons,
only the URI changes. The admin panel adds a confirmation modal before
dispatching the DELETE to guard against accidental deletions.

\begin{infobox}
\textbf{Note on derived data:} Metrics such as race wins, podium counts, and
championship points are \emph{derived} by aggregating race result triples at
query time. They are never stored directly and therefore do not need to be
maintained through UPDATE operations. Deleting a driver removes only the driver
entity; the result triples that reference it are preserved, though they become
orphaned references.
\end{infobox}

\subsubsection{Batch Import, Preview, and Rollback}

The admin panel was extended with a race-results import workflow under
\texttt{/admin-panel/imports/race-results/}. An administrator selects a race,
uploads a CSV file, and receives a dry-run preview before any SPARQL UPDATE is
sent to GraphDB. The preview shows:

\begin{itemize}[leftmargin=1.5em]
  \item new result triples to be inserted;
  \item existing result triples that will be replaced;
  \item unresolved references (driver, constructor, or status IDs not found in the graph);
  \item duplicate or conflicting rows in the uploaded CSV.
\end{itemize}

Only after explicit confirmation is the generated SPARQL batch update executed.
Each confirmed import stores a corresponding rollback record in Django's
relational database, including enough metadata to reverse the exact SPARQL
operation later. The admin dashboard exposes an ``undo latest batch'' style
workflow for these recorded imports.

% =============================================================================
\section{Application Functionalities (UI)}
% =============================================================================

The application is built with Django, rendering server-side HTML templates
and communicating with GraphDB over HTTP. There is no JavaScript framework; all
interactivity is achieved with vanilla JS and CSS custom properties.

\subsection{Public Interface}

\subsubsection{Homepage}

The homepage presents an overview of the knowledge graph with five live stat
tiles (seasons, circuits, drivers, constructors, races) fetched via a single
aggregation query. Below the stats, five navigation cards link to the main
entity sections. A connectivity indicator in the footer shows whether the
GraphDB endpoint is reachable. The page applies a dark motorsport aesthetic
with the F1 red accent colour throughout.

\subsubsection{Drivers Section}

The drivers list page provides:
\begin{itemize}[leftmargin=1.5em]
  \item \textbf{Podium display} --- the top three all-time wins leaders shown as
        styled cards with photo slots for featured drivers.
  \item \textbf{Filterable table} --- server-side filtering by name, nationality,
        and medal type (gold/silver/bronze podiums).
  \item \textbf{Sortable columns} --- clicking any column header re-queries the
        server with the chosen sort dimension.
  \item \textbf{Cards / table toggle} --- client-side view switch persisted in
        \texttt{localStorage}.
  \item \textbf{Career span} --- each driver shows their first and last active
        season (e.g.\ ``2007--2016'').
\end{itemize}

The driver detail page shows biographical info, career statistics (races, wins,
podiums, championships), a horizontal constructor timeline, top circuits, and a
table of the driver's best results.

\subsubsection{Constructors Section}

The constructors list page mirrors the drivers page in structure, providing a
top-three wins podium (with logo image slots for prominent teams), a
nationality filter, server-side column sorting, and a card/table view toggle.
Each card displays the team's nationality, estimated founding year, total race
count, and win tally.

The constructor detail page is divided into four areas:

\begin{itemize}[leftmargin=1.5em]
  \item \textbf{Hero section} --- team name, nationality, and Wikipedia link.
  \item \textbf{Stat tiles} --- total races entered, total wins, total podiums,
        and the number of distinct drivers who raced for the team.
  \item \textbf{Driver roster table} --- every driver who competed for the
        constructor, with their first and last season, race count, wins, and
        podium count. Rows are sorted by wins descending. These figures are
        computed entirely in SPARQL using \texttt{SUM(IF(\ldots))} conditional
        aggregation; no post-processing is required.
  \item \textbf{Top circuits} --- the six circuits where the constructor has
        entered the most races, shown as labelled chips. Useful for identifying
        home-track advantages or high-frequency venues.
\end{itemize}

\subsubsection{Circuits Section}

The circuits list page features a podium of the three circuits that have hosted
the most Grands Prix, each with an image slot for iconic venues (Monaco,
Monza, Silverstone). Below the podium, a card/table grid lists all circuits
with their country, location, and race count badge. Retired circuits are
visually distinguished with a ``Retired'' status chip.

The circuit detail page provides:

\begin{itemize}[leftmargin=1.5em]
  \item \textbf{Hero section} --- circuit name, location, country, and
        geographic coordinates (latitude, longitude, altitude).
  \item \textbf{Stat tiles} --- total races hosted, first and last year of use.
  \item \textbf{Race history table} --- all races held at the circuit listed
        in reverse chronological order, showing year, race name, race winner,
        and winning constructor. Data is fetched with a single SPARQL query
        using \texttt{OPTIONAL} to gracefully handle historic races with missing
        result data.
  \item \textbf{Fastest laps} --- the five all-time fastest individual laps
        recorded at the circuit, sourced from the lap times data
        (filtered to lap duration $> 60{,}000$\,ms to exclude incomplete
        laps) and ordered by ascending milliseconds.
\end{itemize}

\subsubsection{Races Section}

The races list shows a featured card for the most recent race with full podium
display, followed by a filterable list of all races. The race detail page
features:
\begin{itemize}[leftmargin=1.5em]
  \item Hero section with circuit, country, date, and round.
  \item Stat tiles: starters, race laps, fastest lap time and driver.
  \item Podium display (top~3 finishers).
  \item Full results table with grid position, laps completed, race time/gap,
        \textbf{finish status} (e.g.\ ``+1 Lap'', ``Engine'', ``Collision''), and points.
  \item \textbf{Pit stop section} showing per-driver stop count, total pit time,
        and fastest individual stop.
  \item Qualifying results with Q1/Q2/Q3 times.
\end{itemize}

\subsubsection{Seasons Section}

The seasons list page renders one card per season sorted most-recent first.
Each card shows the year prominently, race count, driver count, team count, and
the driver with the most race wins that year. Client-side search and sort
(by year or race count) operate over the rendered cards and a mirrored hidden
table without any server round-trips. The view defaults to cards rather than
table, reflecting the browsing-first use case for seasons.

The season detail page is the richest summary page in the application:

\begin{itemize}[leftmargin=1.5em]
  \item \textbf{Race calendar table} --- every round of the season with round
        number, race name, circuit, country, date, and race winner with their
        constructor. The winner is derived in SPARQL using
        \texttt{OPTIONAL \{ ?res f1:positionOrder 1 \}} so rounds with no
        result data (e.g.\ future races) are still listed.
  \item \textbf{Driver championship standings} --- final standings showing
        position, driver name, total points, and wins. Because the
        \texttt{driver\_standings} table stores cumulative snapshots after
        every round, the query orders by round descending and Python retains
        only the first (latest) row per driver, yielding the end-of-season
        standings efficiently.
  \item \textbf{Constructor championship standings} --- same approach applied
        to \texttt{constructor\_standings}, showing team, points, and wins.
\end{itemize}

\subsubsection{F1 Assistant}

The public interface includes an \texttt{/assistant/} page — accessible via the
last item in the navigation bar — where users can ask natural-language questions
about the Formula 1 knowledge graph. The server uses a Gemini language model
with a hardcoded schema/context prompt describing the main classes, predicates,
and identifier conventions of the RDF graph. The page provides six example
questions as clickable chips to help users discover what the assistant can do.
The workflow is two-stage:

\begin{itemize}[leftmargin=1.5em]
  \item Gemini first generates a read-only \texttt{SELECT} SPARQL query.
  \item The application executes that query against GraphDB and sends the
        resulting bindings back to Gemini, which returns a natural-language answer.
\end{itemize}

The generated SPARQL is never shown to the end user. For safety, only
\texttt{SELECT} queries are accepted; any attempt by the model to generate an
update query is rejected server-side.

\subsection{Admin Panel}

The admin panel is accessible at \texttt{/admin-panel/} and requires login with
a Django staff account. It features:

\begin{itemize}[leftmargin=1.5em]
  \item \textbf{Login page} --- standalone page with the site aesthetic, staff-only.
  \item \textbf{Dashboard} --- live entity counts with quick-add shortcuts.
  \item \textbf{Entity management} for Drivers, Constructors, Circuits, Races,
        and Seasons: list view with inline Edit / Delete actions and Add form.
  \item \textbf{Relation-aware editing} --- entity forms can modify RDF object
        properties such as the team/constructor linked to a driver.
  \item \textbf{Race results CSV import} --- upload, validate, preview, and
        confirm batch result imports for a selected race.
  \item \textbf{Data quality checks} --- dedicated page highlighting races
        without circuits, drivers without constructors, duplicate IDs/labels,
        dangling references, and broken result links.
  \item \textbf{Rollback of batch operations} --- the most recent recorded
        batch import can be reversed from the admin interface.
  \item \textbf{Confirm-delete modal} --- clicking Delete opens a modal
        requiring confirmation before the SPARQL DELETE is executed.
  \item \textbf{Sidebar navigation} --- sticky left nav with active state
        highlighting the current section.
  \item \textbf{Flash messages} --- success and error messages displayed after
        each operation.
\end{itemize}

\subsection{Design System}

The UI is built on a set of CSS custom properties (\texttt{--red},
\texttt{--dark}, \texttt{--panel}, \texttt{--grey}, etc.) that support both
dark and light themes. The theme is persisted in \texttt{localStorage} and
applied before paint to avoid a flash of unstyled content. Typography uses
\emph{Barlow Condensed} (headings) and \emph{Barlow} (body text) from Google
Fonts. The colour palette mirrors the official F1 brand: red (\texttt{\#E10600})
on a near-black background (\texttt{\#15151E}).

% =============================================================================
\section{Conclusions}
% =============================================================================

This project successfully demonstrates the end-to-end lifecycle of a semantic
web application: from raw relational CSV data to an RDF knowledge graph, through
SPARQL queries, to a fully functional web interface.

\subsection{Key Achievements}

\begin{itemize}[leftmargin=1.5em]
  \item \textbf{Full RDF pipeline}: A Python script converts 13 CSV files into
        a consistent N-Triples RDF graph with correct datatype annotations,
        object property links between entities, and \texttt{rdfs:label} triples
        for the main browsable entities.

  \item \textbf{Efficient SPARQL queries}: The anchoring strategy (using unique
        property predicates rather than \texttt{rdf:type}) proved essential for
        acceptable query performance over a graph of millions of triples. Queries
        that naively used \texttt{rdf:type} were orders of magnitude slower.

  \item \textbf{Derived data without denormalisation}: Metrics such as wins,
        podiums, and championship points are computed at query time using SPARQL
        aggregation functions (\texttt{COUNT}, \texttt{SUM(IF(\ldots))},
        \texttt{MIN}, \texttt{MAX}). This keeps the graph normalised and ensures
        consistency.

  \item \textbf{Full CRUD via SPARQL UPDATE}: The admin interface demonstrates
        that GraphDB can serve as the single source of truth for both reads and
        writes, with no relational database involved for domain data.

  \item \textbf{Operations-oriented admin tooling}: Beyond CRUD, the admin now
        supports relation editing, batch race-result imports with dry-run
        previews, rollback of recorded import operations, and graph integrity
        checks.

  \item \textbf{Natural-language query interface}: A Gemini-backed assistant
        translates user questions into safe read-only SPARQL, executes them
        against GraphDB, and returns human-readable answers.

  \item \textbf{Aesthetic and usability}: The application provides a polished
        consumer-grade interface with dark/light theme switching, responsive
        layout, card/table view toggles, client-side search and sort, and
        accessible navigation.
\end{itemize}

\subsection{Challenges and Lessons Learned}

\begin{itemize}[leftmargin=1.5em]
  \item \textbf{xsd:gYear comparison}: Race years stored as \texttt{xsd:gYear}
        cannot be compared directly with integer literals; they must be
        coerced to strings using \texttt{STR()}.

  \item \textbf{Sprint results contamination}: Sprint race results shared the
        \texttt{f1:resultId} anchor and \texttt{f1:race} link with regular
        results, causing duplicate rows. The issue was resolved by filtering
        on the URI pattern of the result entity, although the sprints data was 
        excluded from the graph entirely to prevent future issues.

  \item \textbf{Status join ambiguity}: Joining via \texttt{f1:statusId}
        required care because result entities and status entities both carry this
        property; an additional \texttt{FILTER} on the status entity URI was
        needed to prevent cross-joins.

  \item \textbf{GROUP BY and SAMPLE()}: Season calendar queries must return one
        winner row per race when winners are represented as separate result
        entities. \texttt{SAMPLE()} allowed grouping while retaining an
        arbitrary winner value without requiring a subquery.
\end{itemize}

\subsection{Future Work}

\begin{itemize}[leftmargin=1.5em]
  \item Extend the batch import pattern beyond race results to qualifying,
        pit stops, lap times, and standings.
  \item Implement a public SPARQL explorer endpoint, allowing power users to
        run custom queries against the knowledge graph.
  \item Integrate live data from the current season via the official Formula 1 data feeds or the Kaggle dataset updates.
\end{itemize}

% =============================================================================
\section{Configuration Required to Run the Application}
% =============================================================================

\subsection{Prerequisites}

\begin{itemize}[leftmargin=1.5em]
  \item \textbf{Python 3.11+} with \texttt{pip}
  \item \textbf{GraphDB 10} (free edition) running locally or remotely
\end{itemize}

\subsection{Step 1 --- Clone and Install Dependencies}

\begin{lstlisting}[style=shell]
# From the project root:
python -m venv .venv
source .venv/bin/activate          # Windows: .venv\Scripts\activate
pip install -r requirements.txt
\end{lstlisting}

\subsection{Step 2 --- Configure Environment Variables}

Create a \texttt{.env} file in the project root:

\begin{lstlisting}[style=shell]
DJANGO_SECRET_KEY=your-secret-key-here
DJANGO_DEBUG=True
GRAPHDB_BASE_URL=http://localhost:7200
GRAPHDB_REPOSITORY=ws-formula1
GEMINI_KEY=your-gemini-api-key
# Optional: override the default Gemini model
# GEMINI_MODEL=gemini-2.5-flash
# Optional: if GraphDB requires authentication
# GRAPHDB_USERNAME=admin
# GRAPHDB_PASSWORD=your-password
\end{lstlisting}

\subsection{Step 3 --- Prepare the Knowledge Graph}

\begin{enumerate}[leftmargin=2em]
  \item Place the CSV files from the \textbf{Formula 1 World Championship (1950 - 2024)}
(\url{https://www.kaggle.com/datasets/rohanrao/formula-1-world-championship-1950-2020?select=circuits.csv}) in \texttt{data/raw/}.
  \item Run the conversion script to generate N-Triples:
\begin{lstlisting}[style=shell]
python scripts/csv_to_rdf.py
# Output: data/rdf/f1_dataset.nt
\end{lstlisting}
  \item Open GraphDB Workbench at \texttt{http://localhost:7200}.
  \item Create a repository named \texttt{ws-formula1} with
        \textbf{Ruleset: RDFS+ (Optimized)}.
  \item Go to \textbf{Import $\rightarrow$ Upload RDF files}, select
        \texttt{data/rdf/f1\_dataset.nt}, and import into the named graph
        \texttt{http://example.org/graph/formula1}.
\end{enumerate}

\subsection{Step 4 --- Initialise Django and Create Admin User}

\begin{lstlisting}[style=shell]
python manage.py migrate
python manage.py createsuperuser
# Follow the prompts to set username and password
\end{lstlisting}

\subsection{Step 5 --- Run the Development Server}

\begin{lstlisting}[style=shell]
python manage.py runserver
# Application available at http://127.0.0.1:8000
# Admin panel at http://127.0.0.1:8000/admin-panel/
\end{lstlisting}

\subsection{Summary of URLs}

\begin{center}
\begin{tabular}{@{}ll@{}}
\toprule
\textbf{URL}                    & \textbf{Description} \\
\midrule
\texttt{/}                      & Homepage with live stats \\
\texttt{/assistant/}            & Gemini-backed natural-language assistant \\
\texttt{/drivers/}              & Drivers list \\
\texttt{/constructors/}         & Constructors list \\
\texttt{/circuits/}             & Circuits list \\
\texttt{/races/}                & Races list \\
\texttt{/seasons/}              & Seasons list \\
\texttt{/admin-panel/}          & Admin dashboard (login required) \\
\texttt{/admin-panel/login/}    & Admin login page \\
\texttt{/admin-panel/imports/race-results/} & Race results CSV import \\
\texttt{/admin-panel/data-quality/} & Data quality checks \\
\bottomrule
\end{tabular}
\end{center}

\end{document}
